% Derivation chain figure -- TikZ
% Shows the logical chain from self-modeling to complex QM
% Color code: blue = novel (this paper), gray = published theorem

\begin{figure}[ht]
\centering
\resizebox{\columnwidth}{!}{%
\begin{tikzpicture}[
  % Node styles
  novel/.style={
    rectangle, rounded corners=3pt, draw=blue!70!black, fill=blue!8,
    text width=18em, minimum height=2.4em, align=center,
    font=\small
  },
  published/.style={
    rectangle, rounded corners=3pt, draw=gray!70!black, fill=gray!8,
    text width=18em, minimum height=2.4em, align=center,
    font=\small
  },
  % Arrow style
  arr/.style={->, >=stealth, thick, draw=black!70},
  % Annotation style
  annot/.style={font=\scriptsize\itshape, text=gray!70!black, align=left},
  % Assumption style
  assump/.style={font=\scriptsize, text=red!60!black, align=right}
]

% Nodes (top to bottom)
\node[novel] (L4)
  {\textbf{Self-modeling premise}\\
   (Def.~\ref{def:self-modeling-system}\,(i)--(ii))};

\node[novel, below=0.6cm of L4] (SP)
  {\textbf{Sequential product on} $\boldsymbol{E(V)}$\\
   Corrected product via Peirce feedback};

\node[novel, below=0.6cm of SP] (S17)
  {\textbf{S1--S7 verified}\\
   844+ symbolic tests (\Cref{app:numerics})};

\node[published, below=0.6cm of S17] (EJA)
  {\textbf{Euclidean Jordan algebra}};

\node[novel, below=0.6cm of EJA] (LT)
  {\textbf{Local tomography}\\
   Faithful tracking $\to$ trace form non-degeneracy};

\node[novel, below=0.6cm of LT] (TE)
  {\textbf{Type exclusion}\\
   $\mathbb{R}$, $\mathbb{H}$, spin, Albert excluded};

\node[published, below=0.6cm of TE] (CStar)
  {$\boldsymbol{C^*}$\textbf{-algebra} $\boldsymbol{M_n(\mathbb{C})}$};

\node[published, below=0.6cm of CStar] (INV)
  {\textbf{Involution} $X^* = X^\dagger$\\
   L\"uders product recovered};

% Arrows
\draw[arr] (L4) -- (SP);
\draw[arr] (SP) -- (S17);
\draw[arr] (S17) -- (EJA);
\draw[arr] (EJA) -- (LT);
\draw[arr] (LT) -- (TE);
\draw[arr] (TE) -- (CStar);
\draw[arr] (CStar) -- (INV);

% Right-side annotations (published theorem citations)
\node[annot, right=0.8cm of EJA]
  {van de Wetering, Thm.~1\\
   \cite{vandeWetering2019b}};

\node[annot, right=0.8cm of CStar]
  {van de Wetering, Thm.~3\\
   Barnum--Wilce~\cite{BarnumWilce2014}\\
   Hanche-Olsen~\cite{HancheOlsen1985}};

% Left-side annotations (standing assumptions)
\node[assump, left=0.8cm of LT]
  {Def.~\ref{def:self-modeling-system}\,(iii)--(iv)\\(minimal composite,\\simplicity)};

% Legend
\node[novel, minimum height=1.4em, text width=6em,
      below left=1.2cm and 0.2cm of INV] (leg1)
  {\footnotesize Novel (this paper)};
\node[published, minimum height=1.4em, text width=7em,
      right=0.4cm of leg1] (leg2)
  {\footnotesize Published theorem};

\end{tikzpicture}%
}% end resizebox
\caption{Logical chain of the derivation.  Blue nodes indicate novel
contributions of this paper; gray nodes indicate conclusions drawn from
published theorems (cited).  The four clauses of
Definition~\ref{def:self-modeling-system} enter at the points
indicated.  Every step in the chain is either proved in the main text
or deferred to \Cref{app:proofs}.}
\label{fig:chain}
\end{figure}
